\usepackage{relsize} % for \mathlarger
% Colors tuto: https://www.overleaf.com/learn/latex/Using_colors_in_LaTeX
% \defincolor: {define new color}
%   - rgb: {r, g, b}                            [0, 1]
%   - RGB: {r, g, b}                            [0, 255]
%   - hsb: {hue, saturation, brightness}        [0, 1]
%   - HSB: {hue, saturation, brightness}        [360°, 100%, 100%]
%   - cmy: {cyan, magenta, yellow}              [0, 1]
%   - cmyk: {cyan, magenta, yellow, blacK}      [0, 1]
%   - gray: {gray}                              [0, 1] (low is dark)
%   - Gray: {Gray}                              [0, 15] (low is dark)
%   - HTML: {RRGGBB}                            6 hexa-digits
%   - wave: {wavelength}                        [363, 814] (nm)
% \colorlet: (i.e. mixing colors)
%   - \colorlet{mixedColor}{blue!10!red}        % 10% blue mixed with 90% red

\definecolor{unHighlight}{gray}{0.60}

\newcommand{\bigO}[1]{$\mathcal{O}(#1)$}

\newcommand{\mathtext}[1]{\text{\scriptsize #1}}
% \newcommand{\indexsize}{\footnotesize}
% \renewcommand{\index}[1]{\mathtext{\indexsize {#1}}}
\renewcommand{\index}[1]{\mathtext{{#1}}}
\newcommand{\inmath}[1]{\relax\ifmmode #1 \else $#1$\fi}

\newcommand{\G}[1]{\inmath{G_\index{#1}}}
\newcommand{\N}[1]{\inmath{N_\index{#1}}}
\newcommand{\sk}[1]{\inmath{k_\index{#1}}}
\newcommand{\PK}[1]{\inmath{K_\index{#1}}}
\newcommand{\s}[1]{\inmath{s_\index{#1}}}
\renewcommand{\S}[1]{\inmath{S_\index{#1}}}

\newcommand{\sG}[2] {\relax
    \ifmmode \s{#1}\G{#2} 
    \else $\s{#1}\textcolor{unHighlight}{\G{#2}}$ 
\fi}
\newcommand{\B}[2]{\relax
    \ifmmode \mathlarger{\beta_\index{#1#2}} 
    \else $\mathlarger{\beta_\index{#1\textcolor{unHighlight}{#2}}}$ 
\fi}
\newcommand{\BB}[2]{\relax
    \ifmmode \mathlarger{\beta^\prime_\index{#1#2}} 
    \else $\mathlarger{\beta^\prime_\index{#1\textcolor{unHighlight}{#2}}}$ 
\fi}
\newcommand{\A}[1]{\inmath{\mathlarger{\mathlarger{\alpha_\index{#1}}}}}
\newcommand{\C}[1]{\inmath{\mathlarger{\mathlarger{\gamma_{\index{#1}}}}}}
\newcommand{\Cred}[1]{\inmath{\mathlarger{\mathlarger{\sigma_{\index{#1}}}}}}
\newcommand{\D}{\inmath{\mathlarger{\mathlarger{\mathrm{\Delta}}}}}

\newcommand{\hash}[1]{\inmath{\mathsf{hash}(#1)}}
\newcommand{\sHashed}[1]{\inmath{\sigma_\index{#1}}}
\newcommand{\hashIt}[2]{\inmath{\mathsf{h}^{(#2)}(#1)}}
\newcommand{\Map}[1]{\inmath{\mathsf{Encode}(#1)}}
\newcommand{\InverseMap}[1]{\inmath{\mathsf{Decode}(#1)}}
\newcommand{\InverseMapNoParenthesis}{\inmath{\mathsf{Decode}}}

\newlength{\y}
\newcommand{\width}{9mm}
\newcommand{\height}{5.4mm}

\newcommand{\vgap}{
    \addtolength{\y}{-\height}
    \addtolength{\y}{-1mm}
}
\newcommand{\vGap}{
    \addtolength{\y}{-\height}
    \addtolength{\y}{-3mm}
}
% \newcommand{\hmac}{Eq. \ref{eq:integrity}}
% \newcommand{\hmac}{\footnotesize $\sum$}

% Define TikZ styles
\usepackage{tikz}
\usepackage{tikzpeople}
\usetikzlibrary{shapes.geometric, positioning, calc, patterns, patterns.meta, backgrounds, fit, decorations.pathreplacing}

\tikzset{    
    % shape
    ovale/.style={ellipse, draw, ultra thin, minimum width = #1*\width, minimum height = \height, inner sep = 0pt, anchor = west},
    blockspace/.style={ultra thin, minimum width = #1*\width, minimum height = \height, inner sep = 0pt, anchor = west},
    block/.style={draw, blockspace = #1},
    inblock/.style={block = #1, densely dotted}, % dash pattern=on 1pt off 1pt, 
    zero_pad/.style={block = #1, fill = unHighlight, fill opacity=0.3},
    HMAC/.style={draw, ellipse, minimum width = \width, minimum height = \height, inner sep = 0pt, anchor = center},
    eq/.style={draw, circle, inner sep=1pt, fill opacity=0.2, fill=orange, text opacity=1, text=black, font=\normalsize}, %font=\bfseries\large
    arrow/.style={->, shorten <= 1pt, shorten >= 2pt},
    fleche/.style={-Latex, shorten <= 2pt, shorten >= 3pt},
    line/.style={shorten <= 1pt, shorten >= 1pt},
    dotline/.style={dotted},
    between/.style args={#1 and #2}{at = ($(#1)!0.5!(#2)$)}
}

% Put tikz node in caption
\newcommand\eq[1][black]{ % Do not know why we need [black] to works and \eq instead of {\eq}
    \hspace{-2mm}
    \tikz[baseline=(char.base)]{
        \node [eq] (char) {#1};

    }
    \hspace{-2mm}
}

\newcommand{\blindsign}[1]{ \mathord{\scalebox{1.4}{$\blind{\SIGMA}$}}_{#1} }

\newcommand{\cred}[1]{ \mathord{\scalebox{1.4}{$C$}}_{#1} }


\newcommand{\blind}[1]{\widetilde{#1}}


% \renewcommand{\paragraph}[1]{\medskip\noindent{\textbf{#1.}}}


\newlength{\greekcapheight}

\newcommand{\CapGreek}[1]{%
  \mathord{\mathpalette\CapGreekAux{#1}}%
}

\newcommand{\CapGreekAux}[2]{%
  \begingroup
  \mathsurround=0pt\relax
  \settoheight{\greekcapheight}{$#1 N$}%
  \resizebox{7pt}{\greekcapheight}{$#1#2$}%
  \endgroup
}

\newcommand{\ALPHA}{\CapGreek{\alpha}}
\newcommand{\BETA}{\CapGreek{\beta}}
\newcommand{\GAMMA}{\CapGreek{\gamma}}
\newcommand{\ETA}{\CapGreek{\eta}}
\newcommand{\ZETA}{\CapGreek{\zeta}}
% \newcommand{\NU}{\CapGreek{\nu}}
\newcommand{\KAPPA}{\CapGreek{\kappa}}
\newcommand{\PHI}{\CapGreek{\phi}}
\newcommand{\PSI}{\CapGreek{\psi}}
% \newcommand{\THETA}{\CapGreek{\theta}}
\newcommand{\XI}{\CapGreek{\xi}}
\newcommand{\SIGMA}{\CapGreek{\sigma}}
\newcommand{\RHO}{\CapGreek{\rho}}
% \newcommand{\UPSILON}{\CapGreek{\upsilon}}